%% kapa: exam=1 round=21 year=2010 period=1 form=A subject=civil_law subject_ko=민법 seq=25 number=25
%% source: Q-Net 감정평가사 기출문제 (cst00309) artlSeq=5024181 file=extracted/1008/5024181/2010년 제21회 감정평가사 1차시험 기출문제/1교시.hwp
%% conversion: kordoc 4.24.51 -> tools/gen1.py (auto); review=proofread-sonnet; flags=none
\begin{question}{25}
의사무능력자인 甲은 자신의 부동산에 근저당권을 설정해 주고 금융기관 乙로부터 금전을 대출받아 그 전액을 다시 丙에게 대여하였다. 다음 설명 중 옳지 않은 것은? (다툼이 있으면 판례에 의함)
\begin{choices}
  \item 甲이 의사무능력 상태에서 乙과 체결한 대출거래약정과 근저당권설정계약은 모두 무효이다.
  \item 甲은 받은 이익이 현존하는지 여부에 관계없이 乙에게 대출금 전액을 반환하여야 한다.
  \item 乙은 甲에 대하여 부당이득을 이유로 甲의 丙에 대한 채권의 양도를 구할 수 있다.
  \item 乙로부터 받은 이익이 현존하지 않는다는 점에 대한 증명책임은 甲에게 있다.
  \item 甲이 의사무능력 상태에서 丙에게 한 대여행위도 무효이다.
\end{choices}
\end{question}
\answer{2}
